Glaucoma is a condition that can affect the eyes, causing damage to the ocular nerve and causing the patient's visual field to decrease slowly. If glaucoma is left untreated, the patient can be driven to blindness. Symptoms include:
\begin{itemize}
    \item Gradual loss of side vision
    \item blurred vision
    \item red and swollen eyes
    \item Sensitivity to light
\end{itemize}
The disease is quite common for people over 40, so an annual visit to an eye doctor is recommended to detect the presence of the disease and in a positive case, being able to make treatment suitable for the patient to avoid the advance of the disease. Among the tests performed are fundus exams and Optical coherence tomography.

Given this, using the fundus image and clinical data, machine-learning algorithms support professionals with quick, efficient, and automated diagnoses. In addition, the diagnoses may bring an explicability, giving even more support and resources so that the professional has a bold and more fearless decision about the treatment.

Given the above, this project aims to use machine learning models to detect glaucoma. In addition, the project also presents multiple possibilities for modeling the problem, using different strategies and their benefits. Furthermore, this project explains model performance, its computational efficiency, and the explainability behind the answers of the machine learning models.

